% !TEX root = ~/Code/TN_LOWE_RECONSTRUCTION/main.tex
% This is the frontmatter for the text.Ishould include the title, author, date, and an abstract
\title{Low-Energy Neutrino Event Reconstruction in DUNE: Clustering, Vertexing, and Energy Reconstruction}

% Add affiliation if needed
\newcommand{\ciemat}{\small \textit{Centro de Investigaciones Energéticas, Medioambientales y Tecnológicas (CIEMAT), Madrid, Spain}}

\author{\textbf{Sergio Manthey Corchado}\\
    \ciemat\\
}

\date{\today}
\maketitle
\begin{abstract}
    The \ac{DUNE} aims to measure the two most energetic components of the flux of solar neutrinos ($^{8}\text{B}$ \& $hep$), as well as neutrinos from other astrophysical sources such as nearby supernovae or the diffuse supernovae neutrino background. The sensitivity to the $\nu$--Ar charged current channel, with orders of magnitude higher cross--section than the $\nu$ elastic scattering on elements, together with DUNE's high exposure and great energy resolution, present an exceptional advantage in understanding the characteristics of this neutrino flux. This work presents a detailed study of low--energy charged--current neutrino event reconstruction and energy resolution performance in DUNE, under the presence of significant radiological and external backgrounds. Described are the development and validation of clustering algorithms that identify primary ionization signals and associated adjacent energy depositions to account for spatially correlated activity. For the purpose of solar neutrino analyses, this work requires the additional association of charge and light signals. A dedicated energy calibration strategy is employed to recover the interacting neutrino energy. The impact of various background scenarios on the reconstruction is assessed by reporting on the energy resolution achieved across a range of true neutrino energies. Highlighted are the algorithm's effectiveness in maintaining reconstruction fidelity in the low--energy regime. These results demonstrate the robustness of the reconstruction approach and its potential contribution to the solar--$\nu$ physics program in DUNE.
\end{abstract}
\newpage
\tableofcontents
% \listoffigures
% \listoftables
\newpage